% 来源及取材范围见分题文件；此为整理者转录的正文，不是下载原件。
\begin{definition}[59.67.5]
A field $K$ is called $C_r$ if for every $0<d^r<n$ and every $f\in K[T_1,\ldots,T_n]$ homogeneous of degree $d$, there exist $\alpha=(\alpha_1,\ldots,\alpha_n)$, $\alpha_i\in K$ not all zero, such that $f(\alpha)=0$. Such an $\alpha$ is called a nontrivial solution of $f$.
\end{definition}
\begin{lemma}[59.67.7]
Let $k$ be an algebraically closed field. Let $f_1,\ldots,f_s\in k[T_1,\ldots,T_n]$ be homogeneous polynomials of degree $d_1,\ldots,d_s$ with $d_i>0$. If $s<n$, then $f_1=\ldots=f_s=0$ have a common nontrivial solution.
\end{lemma}
\begin{proof}
This follows from dimension theory, for example in the form of Varieties, Lemma 33.34.2 applied $s-1$ times.
\end{proof}
\begin{theorem}[59.67.10, Tsen's theorem]
The function field of a variety of dimension $r$ over an algebraically closed field $k$ is $C_r$.
\end{theorem}
\begin{proof}
For projective space one can show directly that the field $k(x_1,\ldots,x_r)$ is $C_r$ (exercise).

General case. Without loss of generality, we may assume $X$ to be projective. Let $f\in k(X)[T_1,\ldots,T_n]_d$ with $0<d^r<n$. Say the coefficients of $f$ are in $\Gamma(X,\mathcal O_X(H))$ for some ample $H\subset X$. Let $\boldsymbol\alpha=(\alpha_1,\ldots,\alpha_n)$ with $\alpha_i\in\Gamma(X,\mathcal O_X(eH))$. Then $f(\boldsymbol\alpha)\in\Gamma(X,\mathcal O_X((de+1)H))$. Consider the system of equations $f(\boldsymbol\alpha)=0$. Then by asymptotic Riemann--Roch (Varieties, Proposition 33.45.13) there exists a $c>0$ such that
\begin{itemize}
\item the number of variables is $n\dim_k\Gamma(X,\mathcal O_X(eH))\sim ne^rc$, and
\item the number of equations is $\dim_k\Gamma(X,\mathcal O_X((de+1)H))\sim(de+1)^rc$.
\end{itemize}
Since $n>d^r$, there are more variables than equations. The equations are homogeneous hence there is a solution by Lemma 59.67.7.
\end{proof}
\paragraph{Quoted standard prerequisite: asymptotic Riemann--Roch.}
\begin{proposition}[33.45.13, Tag 0BJ8]
Let $k$ be a field. Let $X$ be a proper scheme over $k$ of dimension $d$. Let $\mathcal L$ be an ample invertible $\mathcal O_X$-module. Then
\[\dim_k\Gamma(X,\mathcal L^{\otimes n})\sim cn^d+l.o.t.\]
where $c=\deg_{\mathcal L}(X)/d!$ is a positive constant.
\end{proposition}
\begin{proof}
This follows from the definitions, Lemma 33.45.9, and the vanishing of higher cohomology in Cohomology of Schemes, Lemma 30.17.1.
\end{proof}
\paragraph{Dimension-theoretic input used in Lemma 59.67.7.}
\begin{lemma}[33.34.2, Tag 0B2Q]
Let $k$ be a field and $n\geq0$. Let $X,Y\subset\mathbf A^n_k$ be closed subsets. Assume that $X$ and $Y$ are equidimensional, $\dim(X)=r$ and $\dim(Y)=s$. Then every irreducible component of $X\cap Y$ has dimension $\geq r+s-n$.
\end{lemma}
\begin{proof}
Consider the closed subscheme $X\times Y\subset\mathbf A^{2n}_k$ where we use coordinates $x_1,\ldots,x_n,y_1,\ldots,y_n$. Then $X\cap Y=X\times Y\cap V(x_1-y_1,\ldots,x_n-y_n)$. Let $t\in X\cap Y\subset X\times Y$ be a closed point. By Lemma 33.20.5 we have $\dim_t(X\times Y)=\dim(X)+\dim(Y)$. Thus $\dim(\mathcal O_{X\times Y,t})=r+s$ by Lemma 33.20.3. By Algebra, Lemma 10.60.13 we conclude that
\[\dim(\mathcal O_{X\cap Y,t})=\dim(\mathcal O_{X\times Y,t}/(x_1-y_1,\ldots,x_n-y_n))\geq r+s-n.\]
This implies the result by Lemma 33.20.3.
\end{proof}
